% Russian manuscript, embedded-inclusion version (October 2026).
% Values in guillemets are filled from the records by fill_embedded.py.
\documentclass[11pt,a4paper]{article}
\usepackage[T2A]{fontenc}
\usepackage[utf8]{inputenc}
\usepackage[english,russian]{babel}
\usepackage[margin=25mm]{geometry}
\usepackage{amsmath,amssymb,graphicx}
\usepackage{placeins}
\usepackage[numbers,sort&compress]{natbib}
\usepackage[hyphens]{url}
\usepackage[hidelinks]{hyperref}
\ifx\pdfmapfile\undefined\else
\pdfmapfile{+cm-super-t1.map}\pdfmapfile{+cm-super-t2a.map}
\pdfmapfile{+cm-super-ts1.map}\fi
\setlength{\emergencystretch}{2em}
\renewcommand{\topfraction}{0.85}
\renewcommand{\textfraction}{0.10}
\renewcommand{\floatpagefraction}{0.75}

\title{Напряжения в алюминиевой матрице вокруг включения Al$_{13}$Fe$_4$
при заданной деформации магнитострикции}
\author{И.\,Е.~Михайловский\thanks{Автор для переписки.
\texttt{ilyamihailovsy@gmail.com}} \and Д.\,Е.~Пшонкин}
\date{Московский политехнический университет, Москва, Россия}

\begin{document}
\maketitle

\begin{abstract}
Методом молекулярной динамики при 0~К рассчитаны напряжения вокруг
сферического включения Al$_{13}$Fe$_4$ радиусом 2{,}5~нм внутри алюминия;
включению задана деформация магнитострикции вдоль магнитного поля.
У включения матрица сжата вдоль поля напряжением «POLE_PEAK»~МПа
и растянута поперёк него напряжением «EQ_PEAK»~МПа, а касательные
напряжения достигают «SHEAR_PEAK»~МПа. Отклик упругий: дислокации не
возникают даже при десятикратной деформации, а дислокация над включением
смещается на «DISL_SHIFT»~\AA{} и возвращается после снятия деформации.
Податливый слой искажённого алюминия у границы снижает атомистическое
поле примерно до «A_PCT»\% решения теории упругости. Если для микронных
частиц это ослабление пренебрежимо (не проверено), то при заданной
деформации упругое напряжение превышает предел текучести сплава в слое
толщиной около радиуса частицы. Наблюдаемый прирост скорости ползучести
при включённом поле потребовал бы активационного объёма втрое меньше
измеренного для сплавов Al--Mg--Si или деформации в 6--7 раз больше
измеренной магнитострикции сплавов Fe--Al. Почему эффект сохраняется
после снятия поля, расчёт не объясняет.
\end{abstract}

\noindent\textbf{Ключевые слова}\quad молекулярная динамика;
магнитострикция; включения; напряжения; алюминиевые сплавы

\section{Введение}

Алюминиевые сплавы системы Al--Mg--Si с железосодержащими включениями
после выдержки в постоянном магнитном поле ползут быстрее, чем без неё
(ползучестью называют медленную деформацию под постоянной нагрузкой),
и этот эффект сохраняется после снятия поля
\cite{Skvortsov2019JAP,Skvortsov2019PSS,Friha2024JMMM,Nikolaev2025SciRep}.
Одно из предложенных объяснений связано с магнитострикцией включений,
то есть с изменением их формы в магнитном поле. Матрица мешает
включению изменить форму, и вокруг него возникают напряжения. Эти
напряжения могут сдвинуть дислокации (линейные дефекты кристалла,
движение которых и даёт пластическую деформацию). В работе \cite{Friha2024JMMM} напряжение
на границе включения оценено в 147~МПа при индукции поля 0{,}7~Тл.
Это одно число: из него не видно, как напряжение
распределено вокруг частицы, как быстро оно спадает с расстоянием и какая
его часть действует на дислокации.

В работе рассчитано поле напряжений вокруг одного включения
Al$_{13}$Fe$_4$ в алюминии. Деформация включения задаётся заранее:
молекулярная динамика рассчитывает движение атомов по межатомному
потенциалу, формуле для энергии их взаимодействия, а такой потенциал
не описывает магнетизм, поэтому магнитострикцию нельзя получить из
самого расчёта. Её величина соответствует оценке 147~МПа
\cite{Friha2024JMMM}. Атомистический расчёт сравнивается с решением
теории упругости, а затем по найденному полю оцениваются сила,
действующая на дислокации, и прирост скорости ползучести.

\section{Модель и методика}

\subsection{Расчётная ячейка}

Расчётной ячейкой служит периодический параллелепипед
алюминия размером $12{,}6\times12{,}4\times12{,}6$~нм. В его центре
находится сферическое включение Al$_{13}$Fe$_4$ радиусом $R=2{,}5$~нм
(рис.~\ref{fig:cell}а). Включение вырезано из кристалла Al$_{13}$Fe$_4$,
решётка которого взята из структурных данных \cite{Liu2024CIF}; атомы алюминия,
оказавшиеся ближе 2{,}2~\AA{} к атомам включения, удалены. В ячейке
118\,861 атом: 114\,374 в матрице и 4\,487 во включении. Включение
со всех сторон окружено алюминием. Граничные условия
периодические по всем трём осям, поэтому расчёт описывает
регулярную решётку включений с шагом 12{,}6~нм вдоль $x$ и $z$ и 12{,}4~нм вдоль $y$; их объёмная доля
равна 3{,}3\%.

Сфера выбрана как простейшая компактная форма; реальная форма включений
в расчёте не воспроизводится. Кроме того, поле
напряжений вокруг сферы спадает с расстоянием $r$ от её центра как
$r^{-3}$ и уже на расстоянии двух-трёх радиусов мало, поэтому для расчёта
достаточно ячейки из $10^5$ атомов. Наконец, для сферы есть точное решение теории
упругости \cite{Eshelby1957,Mura1987}, с которым можно сравнить
атомистический расчёт.

Оси ячейки направлены так: $x\parallel[1\bar10]$, $y\parallel[11\bar2]$,
$z\parallel[111]$. Тогда плоскости $(111)$, по которым в алюминии
скользят дислокации, горизонтальны, а направление скольжения
$[1\bar10]$ совпадает с осью $x$. Взаимодействие атомов описывается
потенциалом модифицированного метода погружённого атома с учётом вторых соседей для системы
Al--Si--Mg--Cu--Fe \cite{Jelinek2012}, расчёты выполнены в LAMMPS
\cite{Thompson2022LAMMPS}. Упругие постоянные алюминия в этом
потенциале $C_{11}=110{,}5$, $C_{12}=60{,}9$, $C_{44}=28{,}4$~ГПа
близки к экспериментальным.

\subsection{Заданная деформация магнитострикции}
\label{sec:strain}

Магнитное поле направлено вдоль $x$. Магнитострикция удлиняет
включение вдоль поля и укорачивает поперёк, не меняя его объёма
(рис.~\ref{fig:cell}б):
\begin{equation}
\varepsilon_{xx}=\varepsilon,\qquad
\varepsilon_{yy}=\varepsilon_{zz}=-\varepsilon/2,\qquad
\varepsilon=1{,}94\cdot10^{-3}.
\label{eq:strain}
\end{equation}
Значение $\varepsilon$ получено делением оценки 147~МПа
из~\cite{Friha2024JMMM} на модуль Юнга 75{,}7~ГПа; модуль алюминия
в используемом потенциале вдоль $[111]$ равен 75{,}9~ГПа.
В~\cite{Friha2024JMMM} оценка сделана с модулем 70~ГПа, и при нём
напряжению 147~МПа отвечает деформация $2{,}1\cdot10^{-3}$. Сама магнитострикция включений Al$_{13}$Fe$_4$ не
измерена, поэтому $\varepsilon$ задаётся как входной параметр расчёта. Все
напряжения, приведённые ниже, пропорциональны $\varepsilon$, и для другой
деформации $\varepsilon'$ их достаточно умножить на $\varepsilon'/0{,}194\%$ (на 1{,}08
для $2{,}1\cdot10^{-3}$).

Каждый атом включения смещается на
$\mathbf u=\boldsymbol\varepsilon(\mathbf r-\mathbf r_c)$, где $\mathbf r$
обозначает положение атома до смещения, $\mathbf r_c$ центр включения,
а $\boldsymbol\varepsilon$ деформацию (\ref{eq:strain}). Затем атомы
включения закрепляются, а атомы алюминия свободно смещаются, пока
силы на них не обратятся в нуль (минимизация энергии методом
сопряжённых градиентов). Так включению задаётся форма, которую,
по предположению, оно принимает в матрице под действием поля. Стрелки на
рис.~\ref{fig:cell}б показывают силы, с которыми включение действует
на алюминий; они взяты из радиального напряжения у границы в решении теории
упругости для одиночной сферы.
На концах вдоль поля включение толкает матрицу наружу, а сверху, снизу
и с боков тянет её внутрь. В реальном включении такую форму поддерживали бы
магнитные силы, если магнитострикция такой величины существует.

Деформация (\ref{eq:strain}) описывает включение внутри матрицы.
Свободная частица того же вещества удлинилась бы в поле сильнее,
потому что матрица ей не мешала бы. Если бы 0{,}194\% была свободной
магнитострикцией упругой частицы, все напряжения в матрице составили бы
долю $A=\beta k/(1+\beta(k-1))$ от рассчитанных. Здесь $k$ обозначает отношение модулей
сдвига включения и алюминия, а коэффициент $\beta=0{,}46$ из решения
теории упругости для сферы зависит только от коэффициента Пуассона
(Дополнительные материалы, раздел~S3). Для включения
в 2--2{,}5 раза жёстче алюминия $A=0{,}63$--$0{,}68$.

Отсчётная ячейка без деформации ($\varepsilon=0$) готовится один раз.
Атомы включения неподвижны с самого начала: в этом потенциале
равновесная решётка Al$_{13}$Fe$_4$ на 8\% больше по объёму, чем
в эксперименте, и свободное включение раздувалось бы. Матрица
прогревается 10~пс при 300~К, охлаждается до 0{,}5~К за 10~пс,
и её энергия минимизируется при 0~К. Затем размер ячейки один раз подстраивается
так, чтобы среднее напряжение в алюминии вдали от включения было
равно нулю. Ячейки с деформацией $+\varepsilon$ и $-\varepsilon$
строятся из отсчётной ячейки смещением атомов включения; при минимизации
энергии в них движутся только атомы матрицы, а размеры ячейки не меняются.
Полусумма откликов на $+\varepsilon$ и $-\varepsilon$ показывает,
насколько отклик отклоняется от линейного (Дополнительные материалы, раздел~S2).

\FloatBarrier
\begin{figure}[!htbp]
\centering
\includegraphics[width=\linewidth]{fig_cell_ru}
\caption{Модель. (а)~Сечение расчётной ячейки плоскостью $xz$ через
центр включения, слой толщиной 1~нм (OVITO \cite{Stukowski2010OVITO}); серые точки: атомы алюминия
матрицы; оранжевые и красные: атомы Al и Fe включения. (б)~Заданная
деформация включения (увеличена для наглядности): штриховая окружность
обозначает исходную форму, заливка соответствует форме с деформацией (\ref{eq:strain}).
Стрелки показывают силы, с которыми включение действует на алюминий
по решению теории упругости для одиночной сферы: наружу
на концах вдоль поля и внутрь сверху и снизу; длина стрелки
пропорциональна силе на единицу площади поверхности; числа показывают эту
силу вдоль и поперёк поля.}
\label{fig:cell}
\end{figure}
\FloatBarrier

\subsection{Как определяются напряжения}

Напряжение на каждом атоме матрицы равно его вириалу, делённому на атомный
объём алюминия 16{,}61~\AA$^3$. Вириал атома складывается из произведений
сил, с которыми на него действуют соседи, на векторы, соединяющие его с ними;
в однородном кристалле сумма вириалов, делённая на объём, равна обычному
напряжению. Рассматривается разность напряжений
ячейки с деформацией и отсчётной ячейки, то есть напряжение, которое
добавляет магнитострикция. Растяжение считается положительным,
сжатие отрицательным. Используются три величины.
Нормальное напряжение вдоль радиуса $\sigma_{rr}$ показывает, как
сильно матрица давит на включение или тянет его; с ним сопоставима
оценка 147~МПа. Эквивалентное напряжение по Мизесу $\sigma_{\text{экв}}$
объединяет все компоненты, кроме всестороннего давления, в одно число,
которое сравнивают с пределом текучести, то есть с напряжением, при котором
сплав начинает деформироваться пластически. Касательное напряжение $\tau=\sigma_{xz}$ действует
в горизонтальной плоскости скольжения вдоль оси $x$ и толкает
краевую дислокацию с вектором Бюргерса вдоль $x$. Вектор Бюргерса задаёт
сдвиг кристалла, который остаётся после прохождения дислокации; у краевой
дислокации он перпендикулярен её линии.

Значения усредняются по атомам в сферических слоях толщиной 2~\AA{}
внутри узких конусов вокруг выбранных направлений: вдоль поля
(отклонение до 10$^\circ$), поперёк поля (до 10$^\circ$ от плоскости
$x=0$) и под углами 25$^\circ$, 45$^\circ$ и 67$^\circ$ к полю
(пояса шириной $\pm7{,}5^\circ$). Расстояние $d=r-R$ отсчитывается
от поверхности включения. Поле симметрично относительно оси $x$
с точностью до слабой упругой анизотропии алюминия (зависимости его
упругости от направления), поэтому для карт
на рис.~\ref{fig:map} напряжения усредняются по поворотам вокруг
этой оси.

\subsection{Решение теории упругости}

Атомистический расчёт сравнивается с двумя решениями теории упругости.
Первое из них точное и относится к одиночной сфере в бесконечной
изотропной матрице с той же деформацией включения
\cite{Eshelby1957,Mura1987,JuSun1999}. Второе получено для той же
периодической ячейки методом Фурье \cite{MoulinecSuquet1998}
с кубическими упругими постоянными алюминия из потенциала. Оно учитывает
соседние включения периодической решётки и неизменный размер ячейки
и поэтому сравнивается с атомистическим расчётом напрямую. Первое
решение показывает, каким было бы поле одиночного включения
в большом зерне.

\section{Результаты}

\subsection{Распределение напряжений}

На рис.~\ref{fig:map} показано, как распределено вокруг включения
добавленное напряжение. Вдоль поля, у концов включения, матрица
сжата: включение раздвигает её. Поперёк поля матрица растянута:
включение отходит от неё внутрь. Касательное напряжение в плоскостях
скольжения равно нулю на осях $x$ и $z$ и максимально в четырёх
наклонных лепестках под углами около 25$^\circ$ и 67$^\circ$ к полю;
в соседних лепестках знаки напряжения противоположны. Атомистическая карта
и карта по теории упругости для той же ячейки совпадают по форме, но
атомистические напряжения примерно на пятую часть меньше
(раздел~\ref{sec:decay}). Ближе «LAYER»~\AA{} к границе различается и форма
поля: теория упругости не учитывает атомного строения границы.

\FloatBarrier
\begin{figure}[!htbp]
\centering
\includegraphics[width=\linewidth]{fig_map_ru}
\caption{Напряжения, добавленные деформацией включения, в плоскости
$xz$ через центр включения: эквивалентное напряжение по Мизесу (сверху)
и касательное напряжение в плоскости скольжения $\tau=\sigma_{xz}$
(снизу). Слева: атомистический расчёт; справа: решение теории
упругости для той же периодической ячейки. Поле направлено вдоль $x$;
окружность показывает поверхность включения, белое кольцо вокруг неё
(3~\AA{} матрицы) не закрашено. Пунктир на левой нижней панели:
направления под 25$^\circ$ и 67$^\circ$ к полю.}
\label{fig:map}
\end{figure}
\FloatBarrier

\subsection{Спад напряжений с расстоянием}
\label{sec:decay}

Профили напряжений вдоль выбранных направлений приведены
на рис.~\ref{fig:decay}. Атомистические значения ниже взяты в слоях
в 5--10~\AA{} от поверхности: ближе к ней в узких конусах вдоль поля всего
15--20 атомов на слой, и значения разбросаны на величину до 25~МПа.
Вдоль поля нормальное напряжение сжимающее, около «MD_POLE_PEAK»~МПа; это
величина того же порядка, что и оценка 147~МПа из \cite{Friha2024JMMM}.
Совпадение ожидаемо, потому что деформация $\varepsilon$ получена из этой
оценки. Поперёк поля нормальное напряжение растягивающее, около
«MD_EQ_PEAK»~МПа. Сбоку от включения, у плоскости $x=0$, велико и
напряжение вдоль поля: удлиняясь, включение тянет за собой прилегающий
к нему алюминий, и у самой поверхности $\sigma_{xx}$ там достигает
«MD_EQ_SXX»~МПа. Касательное напряжение в плоскостях скольжения
в лепестках достигает «MD_SHEAR_PEAK»~МПа.

С удалением от включения все напряжения спадают. На расстоянии одного
радиуса от поверхности ($r=2R$) нормальное напряжение вдоль поля
составляет около 54~МПа, касательное в лепестках около 19~МПа,
на расстоянии двух радиусов ($r=3R$) соответственно около 18 и 7~МПа
(решение для одиночной сферы).

Начиная с $d\approx$~«FIT_D»~\AA{} атомистические напряжения близки
к решению теории упругости для той же ячейки, умноженному на «A_FIT»: после
такого умножения они отличаются от него в среднем на «MED_PCT»\%. Множитель
немного зависит от направления: «A_POLE» вдоль поля, «A_EQ» поперёк поля,
«A_L25» и «A_L67» для касательного напряжения в лепестках 25$^\circ$
и 67$^\circ$. Смещения атомов, которые от способа расчёта напряжений не
зависят, уменьшены так же, а дальше 8~\AA{} от границы напряжение каждого
атома равно упругим постоянным алюминия, умноженным на его деформацию,
с точностью 3\%. Значит, матрица деформируется меньше, чем
в решении теории упругости с идеально сцепленной границей. Разница
накапливается в пределах 5--8~\AA{} от границы: в первом слое алюминия
смещение составляет «KIN_RAD»\% упругого по нормали к поверхности
и «KIN_TAN»\% вдоль неё, в 4--6~\AA{} от границы уже «D46R»\% и «D46T»\%,
а дальше почти не меняется. В этом слое решётка алюминия искажена,
плотность первого слоя вдвое ниже объёмной, и атомы несут меньшее
напряжение при той же деформации (Дополнительные материалы, раздел~S2).

Если податливость этого слоя не зависит от размера частицы, ослабление
обратно пропорционально радиусу, и для частицы радиусом 2{,}5~мкм оно
составило бы около «LOSS_UM»\%. Поэтому для реальных частиц ниже
используется решение теории упругости. Это предположение не проверено:
рассчитаны один радиус включения и одна граница (атомы алюминия ближе
2{,}2~\AA{} к включению удалены) при 0~К. Проскальзывание по границе при
рабочей температуре дало бы ослабление, которое от размера частицы не
зависит.

Решение для ячейки и решение для одиночной сферы различаются: у первого
кубическая упругость алюминия, неизменный размер ячейки и соседние
включения. Касательное напряжение в лепестках у них отличается не больше
чем на 2~МПа при $3\le d\le20$~\AA, а нормальные напряжения в ячейке
больше на 13--17\% уже у поверхности и на 30--44\% при $r=2R$. Объёмная доля
включений в сплаве, 0{,}25\%, в десять с лишним раз меньше, чем в ячейке,
поэтому спад вдали от включения показан по решению для одиночной сферы.

Поле одиночной сферы зависит только от отношения $r/R$. Поэтому
решение теории упругости переносится на включения реального размера: для частицы
диаметром 5~мкм ($R=2{,}5$~мкм) расстояния на рис.~\ref{fig:decay}
нужно умножить на тысячу (верхняя шкала). При заданной деформации
эквивалентное напряжение превышает предел текучести сплава, 40--60~МПа \cite{Friha2024JMMM},
до $r\approx\text{«SHELL_VM»}R$. Наибольшее из касательных напряжений
на двенадцати системах скольжения (система скольжения состоит из плоскости
скольжения и направления скольжения в ней) превышает критическое значение
13--20~МПа до $r\approx\text{«SHELL_SLIP»}R$. Это критическое касательное
напряжение равно пределу текучести, делённому на множитель Тейлора 3{,}06
\cite{StollerZinkle2000}, который учитывает разную ориентацию зёрен
в поликристалле. Для частицы диаметром 5~мкм слой, где эквивалентное
напряжение выше предела текучести, имеет толщину «THICK_VM_UM»~мкм. В этом слое упругое
решение завышает напряжения: в сплаве начнут двигаться уже имеющиеся
дислокации, матрица потечёт пластически, и часть напряжения снимется.
Атомистический расчёт этого не проверяет: его матрица состоит из чистого
алюминия без исходных дислокаций, а новые дислокации в ней при 0~К не
зарождаются даже при касательном напряжении в лепестках «X10_SHEAR»~МПа
(деформация $10\varepsilon$, Дополнительные материалы, раздел~S2).

\FloatBarrier
\begin{figure}[!htbp]
\centering
\includegraphics[width=\linewidth]{fig_decay_ru}
\caption{Добавленные напряжения в зависимости от расстояния $d$
от поверхности включения: (а)~вдоль поля, (б)~поперёк поля,
(в)~касательное напряжение в лепестках 25$^\circ$ и 67$^\circ$.
Точки: атомистический расчёт (среднее по слою толщиной 2~\AA{}
и его погрешность); сплошные линии: решение теории упругости для
той же ячейки; штриховые: решение для одиночной сферы; заштрихован слой
у границы, где форма атомистического поля отличается от упругого
решения. Вдали от границы точки лежат ближе к нулю, чем сплошные линии: атомистическое
поле составляет около «A_PCT»\% решения для той же ячейки. Горизонтальная
полоса показывает предел текучести сплава 40--60~МПа. Верхняя шкала:
расстояние для частицы радиусом 2{,}5~мкм. (г)~Модуль напряжения одиночной
сферы вдали от неё в зависимости от $r/R$ в логарифмическом масштабе; пунктир
показывает наклон $r^{-3}$.}
\label{fig:decay}
\end{figure}
\FloatBarrier

\subsection{Сила, действующая на дислокации}

Касательное напряжение $\tau$ действует на каждую единицу длины
дислокации с вектором Бюргерса $b=2{,}86$~\AA{} с силой $\tau b$,
направленной вдоль плоскости скольжения (сила Пича--Кёлера
\cite{HirthLothe1982}). Дислокация сдвинется, если эта сила больше
силы, которая её удерживает: сопротивления решётки (напряжения Пайерлса,
то есть наименьшего напряжения, при котором прямая дислокация в чистом
кристалле при 0~К переходит в соседнее положение), атомов примесей и других
дислокаций.

На рис.~\ref{fig:glide} показано касательное напряжение поля включения,
усреднённое вдоль прямой, по которой может проходить краевая дислокация.
Прямая параллельна оси $y$ и лежит над включением на высоте $h$ над
его поверхностью. Длина
дислокации равна периоду ячейки вдоль $y$, 12{,}4~нм. Напряжение
меняет знак при замене $x$ на $-x$: по одну сторону от включения оно
толкает дислокацию вперёд, по другую назад. Наибольшее значение, «LINE_PEAK»~МПа,
достигается при $h=$~«LINE_H»~\AA. На коротком участке линии над самим
включением напряжение больше, до «LOCAL_PEAK»~МПа, и этот участок может
прогнуться, даже если среднего напряжения не хватает, чтобы сдвинуть всю
линию. Среднее по линии зависит от расстояния между включениями
в модели (12{,}4~нм вдоль линии) и на сплав напрямую не переносится.

В сплаве Al--Mg--Si дислокации удерживаются атомами Mg и Si
и выделениями, то есть мелкими частицами соединений Mg и Si внутри зёрен; критическое касательное напряжение сплава, полученное
из его предела текучести, равно 13--20~МПа. По решению для одиночной сферы при заданной деформации поле
превышает это значение только в слое у поверхности включения толщиной
«THICK_SLIP» радиуса. Дальше от включения поле меньше этого значения
и само протолкнуть дислокации через препятствия не может; оно лишь
ускоряет или замедляет их термически активированное движение, при котором
препятствия преодолеваются за счёт тепловых колебаний атомов (оценка
для ползучести ниже).

«DISL_TEXT»

\FloatBarrier
\begin{figure}[!htbp]
\centering
\includegraphics[width=0.75\linewidth]{fig_glide_ru}
\caption{Касательное напряжение, усреднённое вдоль прямой краевой
дислокации длиной 12{,}4~нм, которая проходит над включением на высоте
$h$ над его поверхностью, в зависимости от положения дислокации $x$.
Точки: атомистический расчёт; линии: решение теории упругости
для той же ячейки; заливка: протяжённость включения вдоль $x$.
Правая шкала: сила на единицу длины дислокации
$\tau b$.}
\label{fig:glide}
\end{figure}
\FloatBarrier

\subsection{Оценка влияния на ползучесть}

Связь поля напряжений со скоростью ползучести оценивается в модели
термически активированного скольжения
\cite{KocksArgonAshby1975,CaillardMartin2003}. Оценка относится
к ползучести, пока поле включено и напряжения вокруг включений существуют;
наблюдаемый прирост измерен после снятия поля (см. Обсуждение). Дислокация преодолевает
препятствия за счёт тепловых колебаний атомов. Касательное напряжение $\tau$
увеличивает частоту прыжков вперёд в $\exp(V^*\tau/k_{\mathrm B}T)$ раз
и во столько же раз уменьшает частоту прыжков назад. Здесь $V^*$
обозначает активационный объём, то есть меру того, как сильно барьер
зависит от напряжения, а $k_{\mathrm B}T$ обозначает тепловую энергию. Если вокруг включения
дислокации встречают напряжение обоих знаков одинаково часто, средняя
скорость их движения растёт в $\cosh(V^*\tau/k_{\mathrm B}T)$ раз: это
среднее из множителей $\exp(V^*\tau/k_{\mathrm B}T)$ и
$\exp(-V^*\tau/k_{\mathrm B}T)$ для двух знаков напряжения. Превышение
этого множителя над единицей интегрируется по матрице вокруг одного
включения и умножается на число включений в единице объёма $f/V_p$;
так получается относительный прирост скорости ползучести (вывод
приведён в Дополнительных материалах, раздел~S5)
\begin{equation}
H=\frac{f}{V_p}\int_{r>R}
\Bigl\langle\cosh\frac{V^*\tau_s(\mathbf r)}{k_{\mathrm B}T}-1\Bigr\rangle
\,\mathrm dV ,
\label{eq:creep}
\end{equation}
где $f$ и $V_p$ обозначают соответственно объёмную долю включений
и объём одного включения, $\tau_s$ обозначает касательное напряжение
на системе скольжения $s$ (плоскости
скольжения вместе с направлением скольжения в ней), а
угловые скобки означают среднее по двенадцати системам скольжения
и по ориентациям зёрен. В интеграл подставлено поле одиночной сферы по решению теории
упругости (атомистическое ослабление для частиц микронного размера принято
пренебрежимым, раздел~\ref{sec:decay}); оно зависит только от $r/R$,
поэтому $H$ от размера частиц не зависит. Принято $T=300$~К и $f=0{,}246\%$,
что соответствует 0{,}35~мас.\% включений \cite{Friha2024JMMM}.

На рис.~\ref{fig:creep} сплошная кривая показывает амплитуду $\tau_m$,
нужную для прироста скорости ползучести на 25\%, наблюдавшегося после
выдержки в поле \cite{Skvortsov2019JAP}, при разных
$V^*$, если касательное напряжение спадает от поверхности включения как
$\tau_m(R/r)^3$; с ростом активационного объёма эта амплитуда убывает.
Штриховая кривая показывает амплитуду $\tau_m$, при которой такое поле
даёт тот же прирост, что и рассчитанное поле (Дополнительные материалы,
раздел~S5).
Кривые пересекаются при $V^*\approx«VSTAR_GRAINS»\,b^3$. При $V^*=70\,b^3$,
измеренном для сплава Al--Mg--Si \cite{Soula2022JALCOM}, рассчитанное
поле дало бы прирост скорости в $10^8$ раз. Это значение лежит вне области
применимости модели; для наблюдаемых 25\% хватило бы
гораздо меньшей деформации. Поскольку $H$ зависит только от
произведения $V^*\varepsilon$, при $V^*=70\,b^3$ нужна деформация
включения около «EPS_AT_70»\%. Измеренная магнитострикция сплавов Fe--Al
не превышает 0{,}01\% \cite{Hall1959,BormioNunes2012}, то есть
в 6--7 раз меньше этой величины.

\FloatBarrier
\begin{figure}[!htbp]
\centering
\includegraphics[width=0.7\linewidth]{fig_creep_ru}
\caption{Амплитуда касательного напряжения, нужная для прироста
скорости ползучести на 25\% (сплошная линия), и амплитуда,
эквивалентная полю одиночной сферы по теории упругости при деформации 0{,}194\%
(штриховая), в зависимости от активационного объёма $V^*$;
$T=300$~К, $f=0{,}246\%$. Пунктир: $V^*=70\,b^3$ для Al--Mg--Si
\cite{Soula2022JALCOM}.}
\label{fig:creep}
\end{figure}
\FloatBarrier

\section{Обсуждение}

Расчёт отвечает на механическую часть вопроса: какие напряжения
создаёт в алюминии включение заданной формы. Магнитная часть остаётся
открытой. Магнитострикция включений Al$_{13}$Fe$_4$ не измерена,
а для массивного Al$_{13}$Fe$_4$ сильный ферромагнетизм не установлен
\cite{Albedah2015Mossbauer,Avers2025Crystals}. Если реальная деформация
включений близка к магнитострикции сплавов Fe--Al, напряжения будут
в 20 раз меньше рассчитанных, а слой, где они превышают предел текучести,
исчезнет. По решению для одиночной сферы наибольшее касательное напряжение
на системах скольжения у поверхности частицы становится больше критического для сплава при деформации включения
больше «EPS_SLIP_LO»--«EPS_SLIP_HI»\%, а эквивалентное напряжение превышает
предел текучести при деформации больше «EPS_VM_LO»--«EPS_VM_HI»\%. Это
деформация включения внутри матрицы; свободная магнитострикция упругой
частицы для этого должна быть больше в $1/A$ раз (раздел~\ref{sec:strain}).

Форма включения в расчёте задана, поэтому его собственная
упругость в задачу не входит. Межатомный
потенциал хорошо передаёт упругость алюминия, но решётку Al$_{13}$Fe$_4$
завышает, поэтому атомы включения закреплены в экспериментальной решётке.
Расчёт статический (минимизация энергии без теплового движения атомов)
и выполнен при 0~К для включений одного размера (радиус 2{,}5~нм),
образующих регулярную решётку; поле соседних включений учтено решением
теории упругости для той же ячейки. Матрица в расчёте состоит из чистого
алюминия без атомов Mg и Si, выделений и исходных дислокаций, поэтому
предел текучести сплава взят из литературы \cite{Friha2024JMMM},
а пластическое течение у включения не моделируется.

Главное ограничение связано с сохранением эффекта: ползучесть в экспериментах
измеряют после снятия поля, а упругое напряжение магнитострикции
исчезает вместе с полем. Если поле сдвигает дислокации у включений,
необратимым может быть само это перемещение; рассчитанный слой, где
напряжение превышает предел текучести, указывает, где искать такие
перемещения.
В единственном рассчитанном случае, для прямой дислокации при 0~К,
дислокация после снятия деформации вернулась на место; сохраняется ли
такая перестройка дефектов в сплаве и как она влияет на последующую
ползучесть, расчёт не показывает.

\section{Заключение}

Включение Al$_{13}$Fe$_4$ радиусом 2{,}5~нм, которому задана деформация
магнитострикции 0{,}194\% вдоль поля, сжимает алюминиевую матрицу
вдоль поля напряжением около «MD_POLE_PEAK»~МПа и растягивает её поперёк
поля напряжением около «MD_EQ_PEAK»~МПа. Касательное напряжение в плоскостях скольжения
достигает «MD_SHEAR_PEAK»~МПа в наклонных лепестках. В статическом атомистическом расчёте при 0~К
отклик матрицы упругий: дислокаций не возникает даже при деформации в десять раз больше,
а дислокация над включением смещается полем на «DISL_SHIFT»~\AA{}
и возвращается на место, когда деформацию снимают. Атомистическое поле составляет около
«A_PCT»\% решения теории упругости для идеально сцепленного включения;
ослабление создаёт податливый слой искажённого алюминия у границы. Если его
податливость не зависит от размера частицы, для частиц микронного размера
ослабление пренебрежимо мало, но при одном радиусе включения и 0~К это не
проверено. Вдали от включения поле спадает как
$(R/r)^3$; по решению теории упругости при заданной деформации
напряжение превышает предел текучести сплава в слое вокруг
частицы толщиной около её радиуса (для частицы диаметром 5~мкм это
«THICK_VM_UM»~мкм). В этом слое сплав может деформироваться пластически, то есть
необратимым движением уже имеющихся дислокаций; атомистически этот случай
не рассчитан. В модели термически активированного скольжения это поле,
если бы оно действовало во время ползучести,
дало бы наблюдаемый прирост скорости ползучести при активационном объёме
около «VSTAR_GRAINS»$\,b^3$, а при измеренном для сплава $70\,b^3$ для этого хватило бы
деформации «EPS_AT_70»\%,
в 6--7 раз большей измеренной магнитострикции сплавов Fe--Al. Почему эффект
сохраняется после снятия поля, расчёт не объясняет.

\section*{Вклад авторов}

И.\,Е.~Михайловский участвовал в разработке методики
и программного обеспечения, проведении расчётов, анализе, визуализации
и подготовке текста. Д.\,Е.~Пшонкин участвовал в разработке концепции,
руководстве, рецензировании и редактировании текста.

\section*{Доступность данных и кода}

Входные файлы LAMMPS, программы построения ячейки, анализа и рисунков,
отсчётная и деформированные конфигурации с атомными напряжениями
и записи результатов доступны в репозитории
\url{https://github.com/iluuua/magnetostriction-plasticity-bounds}.

\section*{Конфликт интересов}

Авторы заявляют об отсутствии конфликта интересов.

\bibliographystyle{unsrtnat}
\bibliography{references}
\end{document}
